% !TeX root = ../Book4.tex
\begin{figure}[htbp]
	\centering
	\begin{tikzpicture}[font=\small,x=1.05cm,y=1.05cm]
		\fill[GeometryBlueLight] (0,0)--(-1.35,2.1) arc[start angle=180,end angle=360,x radius=1.35,y radius=0.38]--cycle;
		\draw[GeometryBlue,thick] (0,0)--(-1.35,2.1) (0,0)--(1.35,2.1);
		\draw[GeometryBlue,dashed] (1.35,2.1) arc[start angle=0,end angle=180,x radius=1.35,y radius=0.38];
		\draw[GeometryBlue,thick] (-1.35,2.1) arc[start angle=180,end angle=360,x radius=1.35,y radius=0.38];
		\fill[GeometryBlue] (0,0) circle (2pt);
		\node[below] at (0,-0.05) {\(s=0\)：零维层};
		\node[above] at (0,2.58) {\(u^2+v^2=s^2,\ s\ge0\)};
		\node[left] at (-0.78,1.1) {二维层};
		\draw[->,thick] (1.8,1.1)--(3.4,1.1);
		\node[above] at (2.6,1.1) {拓扑同胚};
		\node[below] at (2.6,1.1) {\(w=u+iv\)};
		\fill[GeometryBlueLight] (5,1.1) circle (1.08);
		\draw[GeometryBlue,thick] (5,1.1) circle (1.08);
		\fill[GeometryBlue] (5,1.1) circle (2pt);
		\node[below] at (5,1.02) {\(w=0\)};
		\node[above] at (5,2.58) {平面中的圆盘};
		\node[below] at (5,-0.08) {中心仍须保留约化结构};
	\end{tikzpicture}
	\caption[奇异约化的锥模型]{反向权圆作用的零约化。图中截取了锥的一段；锥顶在拓扑上有圆盘邻域，但其约化光滑函数代数与二维层的辛形式记录了额外的奇异结构。}
	\label{g4:fig:1004:singular-cone}
\end{figure}
